\chapter{From an Isolated Band to a Finite-Range Lattice Hamiltonian}
\index{isolated band!Fourier reduction|textbf}
\index{hopping coefficient|textbf}
\index{finite-range model|textbf}

\label{ch:training-isolated-band-fourier-reduction}

\chapterstatus{pedagogical derivation and synthetic computational example. This
chapter explains the mathematical chain exercised by the controlled
one-dimensional isolated-band calculation, but it reports no M1 numerical
result, material calculation, scientific validation, or uncertainty
quantification.}

\section{Learning goals}

After working through this chapter, a reader should be able to

\begin{enumerate}
  \item construct a one-dimensional Bloch-fiber matrix in plane-wave and
        finite-difference representations;
  \item distinguish a continuum parent, a finite represented reference, and a
        sampled band dispersion;
  \item explain what is retained when one isolated scalar band is selected;
  \item derive the finite reciprocal-to-cell Fourier transform and its inverse;
  \item interpret complete and finite-range hopping coefficients without
        confusing truncation with projection or localization;
  \item derive the Parseval relation for the stated transform convention;
  \item formulate direct hopping recovery as a least-squares problem and state
        the conditions under which it equals Fourier truncation;
  \item keep training and withheld reciprocal samples in distinct roles; and
  \item interpret Hermiticity, bandwidth, curvature, rank, conditioning, and
        imaginary residuals as bounded diagnostics.
\end{enumerate}

Chapter~\ref{ch:training-periodic-geometry} introduced direct and reciprocal
geometry. Chapter~\ref{ch:training-bloch-finite-differences} constructed
Bloch-periodic finite representations, and
Chapter~\ref{ch:training-projection-retained-spaces} distinguished a retained
space from coordinates and later truncations. We now join those ideas in one
complete but deliberately one-dimensional reduction.

The sequence is
\begin{equation}
  \begin{aligned}
    \text{periodic parent}
    &\longrightarrow \text{Bloch fibers}
     \longrightarrow \text{isolated band }E(k),\\
    E(k)
    &\longrightarrow \text{complete hoppings }t_R
     \longrightarrow \text{finite-range model}.
  \end{aligned}
  \label{eq:training-isolated-reduction-sequence}
\end{equation}
Every arrow changes the mathematical object. The purpose of this chapter is to
make those changes visible rather than hiding them behind one plotting routine.

\section{A one-dimensional periodic parent}
\index{cosine potential!pedagogical parent}

Let the direct-lattice period be $a>0$ and define
\begin{equation}
  G=\frac{2\pi}{a}.
  \label{eq:training-isolated-reciprocal-vector}
\end{equation}
Translation by $a$ in real space and translation by $G$ in reciprocal space
supply the fundamental periodicities. Consider the continuum operator
\begin{equation}
  \widehat H
  =-\frac{\hbar^2}{2m}\frac{\mathrm d^2}{\mathrm d x^2}
   +V_0\cos(Gx)
  \quad\text{on }L^2(\mathbb R).
  \label{eq:training-isolated-parent}
\end{equation}
The cosine is a synthetic periodic potential. It is useful because its Fourier
content is transparent; it is not a material potential.

A Bloch state can be written
\begin{equation}
  \psi_k(x)=e^{ikx}u_k(x),
  \qquad
  u_k(x+a)=u_k(x).
  \label{eq:training-isolated-bloch-factorization}
\end{equation}
Substitution into Equation~\eqref{eq:training-isolated-parent} gives the
cell-periodic fiber operator
\begin{equation}
  \widehat H(k)
  =\frac{\hbar^2}{2m}
   \left(-i\frac{\mathrm d}{\mathrm d x}+k\right)^2
   +V_0\cos(Gx).
  \label{eq:training-isolated-fiber}
\end{equation}
The same physics can instead be represented by the original differential
expression on one cell with Bloch boundary conditions. These conventions are
unitarily related, but their matrices are not interchangeable without the
corresponding transformation.

\subsection{Reduced momentum and natural scales}

Define the reduced momentum, reciprocal energy scale, and dimensionless
potential strength by
\begin{equation}
  \kappa=\frac{k}{G},
  \qquad
  E_G=\frac{\hbar^2G^2}{2m},
  \qquad
  \lambda=\frac{V_0}{E_G}.
  \label{eq:training-isolated-scales}
\end{equation}
If $y=Gx$, division by $E_G$ turns the fiber operator into
\begin{equation}
  \widehat h(\kappa)
  =\left(-i\frac{\mathrm d}{\mathrm d y}+\kappa\right)^2
   +\lambda\cos y.
  \label{eq:training-isolated-dimensionless-fiber}
\end{equation}
Thus a statement such as $a=2\pi$, $G=1$, and $E_G=1$ is a dimensionless
representation convention. It does not erase the physical meanings from which
the scales were constructed.

The first Brillouin zone may be represented by the half-open interval
\begin{equation}
  -\frac12\leq\kappa<\frac12.
  \label{eq:training-isolated-first-zone}
\end{equation}
The two endpoints $-1/2$ and $1/2$ differ by one reciprocal-lattice translation
and represent the same boundary point in the reciprocal quotient. A table may
include both endpoints for a symmetry check, but a complete discrete Fourier
mesh should not duplicate them.

\section{Plane-wave representation of a Bloch fiber}
\index{plane-wave representation!one-dimensional fiber}

Because $u_k$ is periodic, expand it in the orthonormal cell basis
\begin{equation}
  \phi_n(x)=\frac{1}{\sqrt a}e^{inGx},
  \qquad n\in\mathbb Z.
  \label{eq:training-isolated-plane-wave-basis}
\end{equation}
Acting on one basis function gives
\begin{equation}
  \left(-i\frac{\mathrm d}{\mathrm d x}+k\right)\phi_n
  =(k+nG)\phi_n.
\end{equation}
The kinetic term is therefore diagonal. The potential identity
\begin{equation}
  V_0\cos(Gx)
  =\frac{V_0}{2}e^{iGx}+\frac{V_0}{2}e^{-iGx}
\end{equation}
shows that multiplication by the cosine changes the reciprocal index by
exactly one. Consequently,
\begin{equation}
  H^{\mathrm{PW}}_{nm}(k)
  =\frac{\hbar^2(k+nG)^2}{2m}\delta_{nm}
   +\frac{V_0}{2}
    \left(\delta_{n,m+1}+\delta_{n,m-1}\right).
  \label{eq:training-isolated-plane-wave-matrix}
\end{equation}
In dimensionless form this becomes
\begin{equation}
  h^{\mathrm{PW}}_{nm}(\kappa)
  =(n+\kappa)^2\delta_{nm}
   +\frac{\lambda}{2}
    \left(\delta_{n,m+1}+\delta_{n,m-1}\right).
  \label{eq:training-isolated-dimensionless-plane-wave-matrix}
\end{equation}

A cutoff $P$ retains $n=-P,\ldots,P$ and gives a
$(2P+1)\times(2P+1)$ Hermitian matrix. Increasing $P$ changes the represented
state space. A cutoff sequence can therefore supply finite-representation
refinement evidence for selected low eigenvalues; it does not make any one
finite matrix identical to the continuum operator.

The nearest-neighbor pattern in reciprocal index deserves careful language. It
comes from the single cosine harmonic. It is not a nearest-neighbor
real-space tight-binding approximation. Reciprocal-index coupling and
real-space hopping are different representations introduced at different
stages.

\section{A finite-difference representation of the same parent}
\index{finite difference!one-dimensional Bloch fiber}

A second representation uses the endpoint-excluded grid
\begin{equation}
  x_j=j\Delta x,
  \qquad
  \Delta x=\frac{a}{N_x},
  \qquad
  j=0,\ldots,N_x-1.
  \label{eq:training-isolated-real-grid}
\end{equation}
Represent the original differential expression on one cell and impose
\begin{equation}
  \psi(x+a)=e^{ika}\psi(x).
  \label{eq:training-isolated-bloch-boundary}
\end{equation}
For the centered second derivative, define
\begin{equation}
  c_h=\frac{\hbar^2}{2m\Delta x^2}.
\end{equation}
The kinetic matrix has diagonal entries $2c_h$, adjacent entries $-c_h$, and
seam entries
\begin{equation}
  T_{0,N_x-1}(k)=-c_h e^{-ika},
  \qquad
  T_{N_x-1,0}(k)=-c_h e^{ika}.
  \label{eq:training-isolated-fd-seam}
\end{equation}
The two seam entries are complex conjugates. Together with the real diagonal
potential samples $V_0\cos(Gx_j)$, this gives a Hermitian finite matrix.

The plane-wave and finite-difference matrices have different bases and usually
different dimensions. Their eigenvalues can be compared as approximations to
the same declared spectral quantities. The matrices themselves cannot be
subtracted until an explicit common-space map has been supplied. Spectral
refinement and represented-operator refinement are therefore different
questions.

\subsection{A finite represented reference is not the continuum limit}

Suppose a high plane-wave cutoff $P_{\mathrm{ref}}$ is declared as the
reference for a finite study. For a candidate cutoff $P$, one may report
\begin{equation}
  \epsilon_P
  =\max_{\kappa\in\mathcal K_{\mathrm{parent}}}
    \max_{1\leq n\leq n_{\mathrm{cmp}}}
    \left|
      E_{n,P}^{\mathrm{PW}}(\kappa)
      -E_{n,P_{\mathrm{ref}}}^{\mathrm{PW}}(\kappa)
    \right|.
  \label{eq:training-isolated-finite-reference-error}
\end{equation}
This is an error relative to the stated finite reference. It is not automatically
$|E_{n,P}-E_{n,\infty}|$. Once discrepancies reach floating-point and
eigensolver scales, they need not decrease monotonically. The accurate report
is the observed finite discrepancy over the tested sequence, unless a separate
argument supplies a continuum bound.

A finite-difference sequence can be compared with the same represented
reference, but its grid and stencil error remains a separate channel from the
plane-wave cutoff error. Agreement of the two sequences is useful precisely
because their discretization mechanisms differ.

\section{What it means to retain an isolated band}
\index{band isolation}

Let $E_n(\kappa)$ be ordered fiber eigenvalues. A band is isolated over a
declared reciprocal domain $\mathcal K$ when it remains separated from the
excluded spectrum. For an interior band, a sampled isolation diagnostic is
\begin{equation}
  \gamma_n^{\mathrm{sampled}}
  =\min_{\kappa\in\mathcal K}
   \min\left\{
     E_n(\kappa)-E_{n-1}(\kappa),
     E_{n+1}(\kappa)-E_n(\kappa)
   \right\}.
  \label{eq:training-isolated-sampled-gap}
\end{equation}
For the lowest band, only the upper gap is present.
A positive value on a finite mesh establishes sampled separation, not a theorem
that the gap is positive at every unsampled point.

At each momentum, four objects must remain distinct:
\begin{enumerate}
  \item the scalar eigenvalue $E_n(\kappa)$;
  \item a normalized eigenvector, whose phase is not unique;
  \item the one-dimensional eigenspace or its projector; and
  \item a choice of phases over momentum, called a band gauge.
\end{enumerate}
For a scalar isolated band, the represented Hamiltonian within the eigenspace is
just the $1\times1$ value $E_n(\kappa)$. Fourier transformation of these scalar
values therefore does not require selecting eigenvector phases. It constructs a
scalar lattice representation of the sampled dispersion. It does not, by
itself, construct Wannier functions, calculate a localization gauge, or retain
a real-space probability density. Those additional objects require eigenvectors
and overlap information~\cite{kohn1959,marzari2012}.

\section{The complete centered reciprocal mesh}
\index{reciprocal mesh!centered uniform}
\index{Born--von Karman representative}

Let $N$ be even. A complete centered half-open mesh is
\begin{equation}
  \kappa_j=-\frac12+\frac{j}{N},
  \qquad j=0,\ldots,N-1.
  \label{eq:training-isolated-centered-mesh}
\end{equation}
It contains $-1/2$ and excludes the equivalent endpoint $1/2$. One convenient
set of cell representatives is
\begin{equation}
  \mathcal R_N
  =\left\{-\frac N2,-\frac N2+1,\ldots,\frac N2-1\right\}.
  \label{eq:training-isolated-representatives}
\end{equation}
The labels in $\mathcal R_N$ are representatives of residue classes modulo
$N$. Thus $R$ and $R+Nq$ produce the same phase at every mesh point. In
particular, the even-mesh endpoint representative $-N/2$ is its own opposite
modulo $N$, because $N/2$ belongs to the same residue class even though it is
not separately stored.

For $R,S\in\mathcal R_N$, finite Fourier orthogonality gives
\begin{equation}
  \sum_{j=0}^{N-1}
  e^{2\pi i(R-S)\kappa_j}
  =N\,\delta_{R,S}^{(N)},
  \label{eq:training-isolated-fourier-orthogonality}
\end{equation}
where $\delta_{R,S}^{(N)}$ equals one when $R$ and $S$ are the same residue
class modulo $N$ and zero otherwise. This identity is the algebraic reason that
the complete transform is invertible.

\section{From sampled energies to complete hoppings}
\index{discrete Fourier transform!band to hopping}

Let $E_j=E(\kappa_j)$ be one real scalar band sampled on the complete mesh.
Define
\begin{equation}
  t_R
  =\frac1N\sum_{j=0}^{N-1}
    e^{-2\pi iR\kappa_j}E_j,
  \qquad R\in\mathcal R_N.
  \label{eq:training-isolated-forward-transform}
\end{equation}
Multiplying by the inverse phases and using
Equation~\eqref{eq:training-isolated-fourier-orthogonality} gives
\begin{equation}
  E_j
  =\sum_{R\in\mathcal R_N}
    e^{2\pi iR\kappa_j}t_R.
  \label{eq:training-isolated-inverse-transform}
\end{equation}
The normalization and signs in these two equations form one convention. Moving
the factor $1/N$ or reversing both signs is possible, but every producer,
interpolator, and verifier must then change consistently~\cite{trefethen2000}.

The coefficient $t_0$ is the cell-diagonal or onsite scalar term. A coefficient
$t_R$ with $R\neq0$ couples cells separated by $R$ direct-lattice periods in
the represented lattice model. The finite mesh produces exactly $N$
independent coefficients. Calling this the complete hopping model means
complete relative to that finite transform, not infinite real-space exactness.

\subsection{Hermiticity and modular opposites}

A scalar lattice Hamiltonian is Hermitian when
\begin{equation}
  t_{\overline{-R}}=t_R^*,
  \label{eq:training-isolated-hopping-hermiticity}
\end{equation}
where $\overline{-R}$ is the selected representative of the opposite residue
class modulo $N$. If the sampled dispersion is real, the complete discrete
transform satisfies this relation up to numerical error. If it is also even in
$\kappa$, the hoppings are real and satisfy
$t_{\overline{-R}}=t_R$.

The modular qualifier matters at the representative boundary. A verifier that
searches only for the literal integer $-R$ can incorrectly report the Nyquist
representative as unpaired. The comparison belongs to the finite periodic
index set, not to an unbounded integer list.

\subsection{A four-point transform by hand}

Consider the already finite scalar model
\begin{equation}
  E(\kappa)=\varepsilon_0+2t\cos(2\pi\kappa).
  \label{eq:training-isolated-hand-band}
\end{equation}
For $N=4$, the mesh is
$\{-1/2,-1/4,0,1/4\}$ and the sampled values are
\begin{equation}
  \varepsilon_0-2t,
  \quad\varepsilon_0,
  \quad\varepsilon_0+2t,
  \quad\varepsilon_0.
\end{equation}
Substitution into Equation~\eqref{eq:training-isolated-forward-transform}
gives
\begin{equation}
  t_0=\varepsilon_0,
  \qquad
  t_{-1}=t_1=t,
  \qquad
  t_{-2}=0.
\end{equation}
The inverse transform recovers all four samples. This example makes the meaning
of ``range one'' literal: only the onsite and adjacent-cell coefficients are
nonzero. A general sampled band need not terminate at any finite range.

\section{Finite-range truncation}
\index{hopping truncation}

For an integer $R_{\max}<N/2$, retain the symmetric representative set
\begin{equation}
  \mathcal R_{R_{\max}}
  =\{R\in\mathcal R_N:|R|\leq R_{\max}\}.
  \label{eq:training-isolated-range-set}
\end{equation}
The Fourier-mediated finite-range model is
\begin{equation}
  E_{R_{\max}}^{\mathrm{med}}(\kappa)
  =\sum_{R\in\mathcal R_{R_{\max}}}
    e^{2\pi iR\kappa}t_R.
  \label{eq:training-isolated-truncated-band}
\end{equation}
This operation removes hopping coefficients. It does not reduce the dimension
of the already selected scalar band, select a new eigenspace, or choose a
Wannier gauge. It is a model-class approximation performed after band
retention.

The omitted-coefficient norm
\begin{equation}
  \eta_{R_{\max}}
  =\left(
    \sum_{R\notin\mathcal R_{R_{\max}}}|t_R|^2
   \right)^{1/2}
  \label{eq:training-isolated-omitted-norm}
\end{equation}
is a representation diagnostic with energy units. It is not a probability, an
energy fraction, or a material uncertainty.

\section{Parseval's identity for the truncation residual}
\index{Parseval identity!finite hopping transform}

On the training mesh define
\begin{equation}
  r_j
  =E_{R_{\max}}^{\mathrm{med}}(\kappa_j)-E_j
  =-\sum_{R\notin\mathcal R_{R_{\max}}}
    e^{2\pi iR\kappa_j}t_R.
  \label{eq:training-isolated-training-residual}
\end{equation}
Then
\begin{align}
  \sum_{j=0}^{N-1}|r_j|^2
  &=\sum_j
    \sum_{R\notin\mathcal R_{R_{\max}}}
    \sum_{S\notin\mathcal R_{R_{\max}}}
    e^{2\pi i(R-S)\kappa_j}t_Rt_S^*\\
  &=N\sum_{R\notin\mathcal R_{R_{\max}}}|t_R|^2.
  \label{eq:training-isolated-parseval}
\end{align}
The second line follows from finite Fourier orthogonality. This is an exact
identity for the stated complete mesh, representatives, and normalization.

Parseval equates a squared residual summed over the training mesh with a
squared omitted-coefficient norm. It does not state that the maximum error is
the same quantity. It also does not determine the error at a new reciprocal
coordinate, where finite Fourier interpolation rather than discrete
orthogonality controls the comparison.

\section{The direct least-squares route}
\index{least squares!direct hopping fit}
\index{route comparison!Fourier mediated}

Instead of first calculating all $N$ coefficients, one can fit a chosen
finite-range class directly. Let the columns be indexed by
$R\in\mathcal R_{R_{\max}}$ and define
\begin{equation}
  A_{jR}=e^{2\pi iR\kappa_j}.
  \label{eq:training-isolated-design-matrix}
\end{equation}
For nonnegative declared weights $w_j$, the direct route solves
\begin{equation}
  \min_{\boldsymbol\tau}
  \sum_{j=0}^{N-1}w_j
  \left|
    E_j-\sum_R A_{jR}\tau_R
  \right|^2.
  \label{eq:training-isolated-weighted-fit}
\end{equation}
The unknown vector $\boldsymbol\tau$ contains the directly fitted hoppings. A
stable numerical implementation normally uses a QR or singular-value method
rather than explicitly forming inverse normal equations~\cite{golubVanLoan2013}.

With the complete uniform mesh, equal weights, and distinct representatives,
Fourier orthogonality gives
\begin{equation}
  A^{\dagger}A=NI.
  \label{eq:training-isolated-orthogonal-design}
\end{equation}
The least-squares solution is therefore
\begin{equation}
  \boldsymbol\tau
  =\frac1N A^{\dagger}\mathbf E,
  \label{eq:training-isolated-fit-fourier-solution}
\end{equation}
which selects the same coefficients as the complete transform and then discards
the unrepresented columns. Under these hypotheses, direct fitting and
Fourier-mediated truncation are the same orthogonal projection.

This equivalence is conditional, not a general property of two workflows. It
can be lost by changing
\begin{itemize}
  \item the reciprocal coordinates or Fourier signs;
  \item the weights;
  \item the fitting domain;
  \item the retained representatives;
  \item Hermiticity or parameter-sharing constraints;
  \item the fitted observable; or
  \item, for a multiband operator, the aligned frame.
\end{itemize}
When one of these changes, a nonzero route defect may correctly report two
different approximation problems rather than a software failure.

\subsection{Rank, identifiability, and conditioning}

The design rank is the dimension of the coefficient combinations visible on
the training coordinates. If
\begin{equation}
  \operatorname{rank}(A)<|\mathcal R_{R_{\max}}|,
\end{equation}
then multiple coefficient vectors produce the same fitted samples and the
individual hoppings are not identified by that design. Full column rank gives a
unique least-squares minimizer, but uniqueness does not guarantee numerical
stability.

If $\sigma_{\max}$ and $\sigma_{\min}$ are the largest and smallest nonzero
singular values, the two-norm condition number is
\begin{equation}
  \kappa_2(A)=\frac{\sigma_{\max}}{\sigma_{\min}}.
  \label{eq:training-isolated-design-condition}
\end{equation}
A large condition number warns that small perturbations in samples may produce
large coefficient changes. On the matched complete uniform design above, the
columns are orthogonal and the condition number is one. That favorable value is
a property of this design, not of every hopping fit.

\section{Training and withheld reciprocal coordinates}
\index{training mesh!complete Fourier}
\index{withheld mesh!staggered reciprocal}

The complete training mesh has a special algebraic role: it defines the
finite Fourier transform. Reconstructing those same $N$ values with all $N$
coefficients verifies the transform pair, but it is interpolation of the data
that defined the coefficients. It is not evidence about unsampled coordinates.

A separate uniform withheld mesh can be staggered so that it shares no training
coordinates. For training extent $N$, withheld extent $M$, and
\begin{equation}
  \eta=\frac{1}{N+1},
\end{equation}
define
\begin{equation}
  \kappa_i^{\mathrm W}
  =-\frac12+\frac{i+\eta}{M},
  \qquad i=0,\ldots,M-1.
  \label{eq:training-isolated-withheld-mesh}
\end{equation}
To see why the meshes are disjoint, suppose a training and withheld coordinate
were equal. Cancelling $-1/2$ would give
\begin{equation}
  \frac{j}{N}=\frac{i+1/(N+1)}{M}.
\end{equation}
Multiplication by $MN(N+1)$ would require
\begin{equation}
  jM(N+1)=iN(N+1)+N.
\end{equation}
The left side is divisible by $N+1$, whereas the right side has remainder $N$
and is not. The assumed equality is impossible.

The withheld values are evaluated only after the coefficients, tested ranges,
and tolerances have been frozen. They may reveal interpolation behavior between
training points. They must not update the fit, choose a range, alter a threshold,
or become a new training set while retaining the label ``withheld.'' In this
synthetic setting, withheld evaluation is still numerical evidence relative to
the selected parent representation; it is not scientific validation.

\section{Band-shape and structural diagnostics}
\index{bandwidth!scalar hopping model}
\index{curvature!reduced momentum}

A diagnostic has meaning only after its domain and coordinate are declared. On
a finite comparison set $\mathcal K_{\mathrm{cmp}}$, the sampled bandwidth is
\begin{equation}
  W_{\mathrm{sampled}}
  =\max_{\kappa\in\mathcal K_{\mathrm{cmp}}}E(\kappa)
   -\min_{\kappa\in\mathcal K_{\mathrm{cmp}}}E(\kappa).
  \label{eq:training-isolated-bandwidth}
\end{equation}
It need not equal the exact continuous bandwidth unless the extrema are known
to lie in the sampled set.

For the trigonometric polynomial
\begin{equation}
  E_R(\kappa)=\sum_R e^{2\pi iR\kappa}t_R,
\end{equation}
differentiation is analytical:
\begin{equation}
  \frac{\mathrm d^2E_R}{\mathrm d\kappa^2}(0)
  =-(2\pi)^2\sum_R R^2t_R.
  \label{eq:training-isolated-reduced-curvature}
\end{equation}
Because $\kappa=k/G$,
\begin{equation}
  \frac{\mathrm d^2E_R}{\mathrm d k^2}
  =\frac{1}{G^2}
   \frac{\mathrm d^2E_R}{\mathrm d\kappa^2}.
  \label{eq:training-isolated-physical-curvature}
\end{equation}
Thus reduced-coordinate curvature and physical-wave-vector curvature have
different units. Neither becomes a material effective mass until the extremum,
physical scales, band identity, and effective-mass convention are supplied.

Numerical interpolation may also produce small imaginary components even when
the target band is real. A scalar diagnostic can report
\begin{equation}
  \epsilon_{\mathrm{imag}}
  =\max\left\{
    \max_{\kappa\in\mathcal K_{\mathrm{cmp}}}
      |\operatorname{Im}E_R(\kappa)|,
    \left|
      \operatorname{Im}
      \frac{\mathrm d^2E_R}{\mathrm d\kappa^2}(0)
    \right|
  \right\}.
  \label{eq:training-isolated-imaginary-residual}
\end{equation}
This checks consistency with a real scalar observable. It does not replace the
coefficient-level Hermiticity check in
Equation~\eqref{eq:training-isolated-hopping-hermiticity}.

Maximum error and root-mean-square error must also be named separately. For
residuals $r_i$,
\begin{equation}
  \epsilon_{\max}=\max_i|r_i|,
  \qquad
  \epsilon_{\mathrm{RMS}}
  =\left(\frac1M\sum_i|r_i|^2\right)^{1/2}.
  \label{eq:training-isolated-max-rms}
\end{equation}
They answer different questions and generally have different values. A report
must not transfer a number from one metric to the other because both are called
``band error.''

\section{A transparent computational example}

The following synthetic example implements the transform, truncation, direct
fit, and withheld comparison directly. It is intentionally smaller than a
research calculation so that every array has an evident meaning.

\begin{pythonexample}{transforming a sampled scalar band into hoppings}
import numpy as np

# Use an even complete mesh on the half-open reduced Brillouin zone.
N = 8
training_kappa = -0.5 + np.arange(N, dtype=np.float64) / N
representatives = np.arange(-N // 2, N // 2, dtype=np.int64)

# Define a synthetic scalar band with range-two Fourier content.
epsilon_0 = 0.7
t_1 = -0.12
t_2 = 0.015
training_energy = (
    epsilon_0
    + 2.0 * t_1 * np.cos(2.0 * np.pi * training_kappa)
    + 2.0 * t_2 * np.cos(4.0 * np.pi * training_kappa)
)

# Apply the declared forward transform one representative at a time.
hoppings = np.array(
    [
        np.mean(
            training_energy
            * np.exp(-2j * np.pi * R * training_kappa)
        )
        for R in representatives
    ],
    dtype=np.complex128,
)

# The complete inverse transform must reproduce every training value.
complete_design = np.exp(
    2j * np.pi * np.outer(training_kappa, representatives)
)
reconstructed = complete_design @ hoppings
np.testing.assert_allclose(reconstructed.real, training_energy, atol=1e-14)
np.testing.assert_allclose(reconstructed.imag, 0.0, atol=1e-14)

# Retain the symmetric range-one model and evaluate its training residual.
range_mask = np.abs(representatives) <= 1
range_representatives = representatives[range_mask]
mediated_hoppings = hoppings[range_mask]
range_design = np.exp(
    2j * np.pi * np.outer(training_kappa, range_representatives)
)
mediated_training = range_design @ mediated_hoppings

# Equal-weight direct least squares on the same complete mesh is equivalent.
direct_hoppings, _, rank, singular_values = np.linalg.lstsq(
    range_design,
    training_energy.astype(np.complex128),
    rcond=None,
)
np.testing.assert_allclose(direct_hoppings, mediated_hoppings, atol=1e-14)
assert rank == range_representatives.size
condition_number = singular_values[0] / singular_values[-1]

# Use a frozen staggered mesh only for evaluation, never for the fit.
M = 17
offset = 1.0 / (N + 1)
withheld_kappa = -0.5 + (np.arange(M) + offset) / M
withheld_target = (
    epsilon_0
    + 2.0 * t_1 * np.cos(2.0 * np.pi * withheld_kappa)
    + 2.0 * t_2 * np.cos(4.0 * np.pi * withheld_kappa)
)
withheld_design = np.exp(
    2j * np.pi * np.outer(withheld_kappa, range_representatives)
)
withheld_candidate = withheld_design @ mediated_hoppings
withheld_maximum_error = np.max(
    np.abs(withheld_candidate.real - withheld_target)
)

print("direct design rank:", rank)
print("direct design condition number:", condition_number)
print("withheld maximum error:", withheld_maximum_error)
\end{pythonexample}

The example's range-one error is intentional because the target contains a
range-two term. Exact complete-mesh reconstruction checks the transform. The
nonzero withheld error checks the deliberately truncated candidate. Neither
observation is a material result.

\documentcallout{Notebook to do}{Create
\path{examples/tutorials/research-monograph/foundations/isolated_band_fourier_reduction.ipynb}
only after this lecture contract is stable. Include explicit imports,
explanatory inline comments, a hand-checkable transform, modular Hermiticity,
Parseval, equal-weight route equivalence, a deliberately changed-weight route,
and disjoint withheld evaluation. Keep all examples synthetic and do not create
or execute the notebook as part of the present manuscript change.}

\section{A staged isolated-band laboratory}
\index{laboratory!isolated-band Fourier reduction}

A useful laboratory should preserve the sequence of objects rather than jumping
directly from a potential to a final error curve.

\begin{enumerate}
  \item Declare $a$, $G$, $E_G$, $\lambda$, the reciprocal interval, and the
        energy zero.
  \item Construct plane-wave fiber matrices and verify their Hermiticity and
        reciprocal-index coupling pattern.
  \item Construct the independent finite-difference fibers, including the
        conjugate seam phases.
  \item Compare a frozen set of low eigenvalues against a declared finite
        reference while retaining the two discretization channels separately.
  \item Check the sampled isolation gap before assigning a scalar band role.
  \item Freeze the complete training mesh, representative convention, tested
        ranges, withheld rule, metrics, and numerical tolerances.
  \item Calculate the complete transform and verify inverse reconstruction and
        modular hopping Hermiticity.
  \item For each range, retain the omitted norm, training and withheld errors,
        Parseval residual, direct-fit rank and conditioning, route defects, and
        band-shape diagnostics.
  \item Interpret numerical-consistency tolerances separately from any later
        model-selection or scientific-acceptance threshold.
\end{enumerate}

If one stage fails, the workflow should stop or retain the failure explicitly.
For example, a poorly isolated band is not repaired by making the hopping range
larger, and a rank-deficient direct design is not repaired by reporting only its
training residual.

\section{Common mistakes}

\begin{enumerate}
  \item Including both reciprocal-zone endpoints in a complete Fourier mesh.
  \item Calling the reciprocal-index tridiagonal plane-wave matrix a
        real-space nearest-neighbor model.
  \item Treating a high finite cutoff as the exact continuum spectrum.
  \item Assuming a positive sampled gap proves isolation between sample points.
  \item Claiming that scalar energy samples construct a Wannier gauge or
        localized orbital.
  \item Ignoring modular opposites at the even-mesh Nyquist representative.
  \item Calling hopping truncation a projection onto fewer quantum states.
  \item Applying Parseval's training-mesh identity directly to a withheld
        maximum error.
  \item Assuming direct fitting and mediated truncation agree after changing
        weights, coordinates, constraints, or model classes.
  \item Using withheld points to choose the hopping range and then continuing
        to call them withheld.
  \item Reporting reduced-coordinate curvature as a physical effective mass.
  \item Treating a binary64 numerical tolerance as scientific uncertainty.
\end{enumerate}

\section{Exercises}

\begin{enumerate}
  \item Starting from Equation~\eqref{eq:training-isolated-fiber}, derive every
        term in Equation~\eqref{eq:training-isolated-plane-wave-matrix}.
  \item Construct the $P=1$ plane-wave matrix explicitly for arbitrary
        $\kappa$ and $\lambda$, and verify its Hermiticity.
  \item Derive the two finite-difference seam entries in
        Equation~\eqref{eq:training-isolated-fd-seam} from the Bloch boundary
        condition.
  \item Prove Equation~\eqref{eq:training-isolated-fourier-orthogonality} by
        summing a finite geometric series.
  \item Complete every sum in the four-point example and explain why the
        representative $-2$ is its own modular opposite.
  \item Derive Equation~\eqref{eq:training-isolated-parseval} without invoking
        an infinite Fourier-series theorem.
  \item Use Equation~\eqref{eq:training-isolated-orthogonal-design} to derive
        the direct-fit solution in
        Equation~\eqref{eq:training-isolated-fit-fourier-solution}.
  \item Change one weight in the Python example and measure the resulting
        direct-versus-mediated coefficient defect.
  \item Prove the withheld-mesh disjointness argument for arbitrary positive
        integers $N$ and $M$.
  \item Convert a reduced curvature to physical-$k$ curvature and list the
        additional assumptions required to call it an effective mass.
\end{enumerate}

\section{Evidence boundary and next use}

The derivations establish finite-transform identities under their stated
conventions. The Python example checks those identities for synthetic data.
Neither establishes that a particular continuum calculation is converged, that
a finite hopping range is adequate for an application, or that a physical
material is described correctly.

Appendix~\ref{app:one-dimensional-reduction} applies this chain to retained
controlled calculations. Its numerical outcomes, historical localization
channels, and provenance remain owned by their calculation packages. The
present chapter supplies the explanatory mathematics and does not inherit their
evidence status. Chapter~\ref{ch:training-composite-band-gauge-alignment} next
replaces the scalar retained space by a composite frame and derives the missing
Bloch--Wannier, transport, alignment, and block-locality relations. The final
Part~I chapter then explains how to classify claims, references, observations,
and dispositions.
